% Enunciado: Seção 5, pontos obrigatórios da discussão
% Tamanho sugerido: ~1 coluna

\subsection{Ótimos locais}
\label{subsec:otimos-locais}
% - Quantos ótimos locais distintos o greedy descent encontrou na
%   instância pequena ao longo dos restarts (somando as 30 execuções).
% - Em quantas execuções cada algoritmo alcançou h* = 1875.
% - Relacionar com ótimo local x ótimo global visto em aula.
% - Opcional (anexo): histograma dos valores de h dos ótimos locais.

\begin{table}[H]
	\centering
	\caption{Instância pequena: ótimos locais e execuções que atingem $h^\star$.}
	\label{tab:otimos-locais}
	\begin{tabular}{lcc}
		\toprule
		Algoritmo & Ótimos locais distintos & Exec. com $h^\star$ (de 30) \\
		\midrule
		\linhasresultado{tab-otimos-locais}
		\bottomrule
	\end{tabular}
\end{table}

\subsection{Exploração $\times$ intensificação}
\label{subsec:exploracao}
% - Para cada algoritmo, qual mecanismo explora e qual intensifica.
% - Onde os dados mostram equilíbrio bem/mal ajustado, apontando figuras:
%   restarts (Tabela~\ref{tab:gd-passos}), taxa de aceitação
%   (Figura~\ref{fig:sa-aceitacao}), diversidade
%   (Figura~\ref{fig:ag-diversidade}) e distância melhor x média do AG
%   (Figura~\ref{fig:convergencia}).

\begin{table}[H]
	\centering
	\caption{Mecanismos de exploração e intensificação de cada algoritmo.}
	\label{tab:exploracao}
	\begin{tabular}{lp{0.33\columnwidth}p{0.33\columnwidth}}
		\toprule
		Algoritmo & Exploração & Intensificação \\
		\midrule
		GD+RR     & --         & --             \\
		SA        & --         & --             \\
		AG        & --         & --             \\
		\bottomrule
	\end{tabular}
\end{table}

\subsection{Efeito da escala}
\label{subsec:escala}
% - O ranking dos três algoritmos muda da instância pequena para a grande?
% - Se sim: hipótese ligada ao custo de varrer a vizinhança (O(n^2)
%   sucessores por passo) e ao tamanho do espaço de busca, sustentada com
%   os dados (Tabelas~\ref{tab:principal}, \ref{tab:gd-passos} e
%   \ref{tab:tempo}).
% - Se não: mostrar com os p-valores que a estabilidade do ranking é
%   estatisticamente sustentada.
% Sem figura própria; reaproveitar as da Seção~\ref{subsec:comparacao}.

\subsection{Quando o algoritmo genético compensa?}
\label{subsec:quando-ag}
% - Existe instância em que o AG supera as duas buscas locais com
%   significância estatística? Citar Tabela~\ref{tab:pvalores}.
% - Hipótese em termos da codificação: a recombinação consegue combinar bons
%   trechos de sequência de dois pais num filho melhor?
% - Experimento adicional que testa a hipótese: desenho (ex.: AG com vs.
%   sem cruzamento, OX vs. PMX, ou taxa de filhos melhores que ambos os
%   pais ao longo das gerações), resultado e p-valor.

\begin{table}[H]
	\centering
	\caption{Experimento adicional: efeito da recombinação.}
	\label{tab:experimento-ag}
	\begin{tabular}{llcc}
		\toprule
		Instância & Configuração & Melhor $h$ (média $\pm$ dp) & p-valor \\
		\midrule
		\linhasresultado{tab-experimento-ag}
		\bottomrule
	\end{tabular}
\end{table}

% Opcional, se o experimento for sobre a dinâmica ao longo das gerações:
% figura com a curva medida (ex.: fração de filhos melhores que os pais
% x geração). Arquivo sugerido: results/ag\_experimento.png
